\documentclass[11pt,a4paper]{ctexart}
\usepackage[margin=2.5cm]{geometry}
\usepackage{graphicx}
\usepackage{booktabs}
\usepackage{amsmath}
\usepackage{hyperref}
\title{面向细粒度情绪识别的多标签文本分类方法对比研究}
\author{课程项目组（待填写）}
\date{\today}
\begin{document}
\maketitle
\begin{abstract}
本报告研究基于 GoEmotions 数据集的英文细粒度多标签情绪分类。\textbf{此处为结构占位，待正式实验完成后填写定量结果与结论。}
\end{abstract}
\section{引言与研究问题}
RQ1：不同统计学习、神经网络和预训练 Transformer 方法的多标签情绪识别表现。RQ2：更强的预训练 Encoder 是否改善细粒度分类。
\section{数据集与预处理}
GoEmotions simplified；官方 train/dev/test；28 个标签；标签频次与不平衡分析；防止数据泄漏。
\section{方法}
\subsection{TF-IDF 与 Linear SVM（E1）}
\subsection{TextCNN（E2）}
\subsection{BiLSTM（E3）}
\subsection{预训练 Transformer（E4）}
\subsection{Word2Vec 词嵌入消融（E5）}
\section{实验设置与验证协议}
统一拆分、训练预算、阈值选择、Micro-F1、Macro-F1、随机种子和 Run 记录。
\section{实验结果与分析}
\textbf{待自动从冻结的 Run 生成表格与图片；不得填入模拟成绩。}
\section{系统实现与演示}
Monorepo 架构、推理 API、React 文本分类与情绪得分可视化。
\section{局限性与结论}
标注主观性、Reddit 领域偏差、类别不均衡与跨领域推广。
\begin{thebibliography}{9}
\bibitem{goemotions} Demszky et al. GoEmotions: A Dataset of Fine-Grained Emotions. ACL, 2020.
\end{thebibliography}
\end{document}
